% Nouveau EGA V (cote 25): archive cote 25, site page 82 (PDF page 83).
% typescript. Draft compared: none.
% AI-transcribed by Claude Opus 5.5 (claude-opus-5-5, adaptive thinking, effort high), Message Batches API, batch ega25b (msgbatch_01A7n5YuXubUPNK1zsepi8RG), 2026-11-13; not yet checked by eye.
\scanpage{25/82}
\mspage{\hand{77}}

de base général. Signalons la suivante (sans démonstration): Soit $f:X\to S$
un morphisme projectif et plat, \xout{\unsure{xxxSxxixixquxxxX}} à fibres géométriques
\xout{\unsure{géomé}} irréductibles et $(R_2)$ , $S'$ un sous-\inserted[typed]{pré}schéma de $X$ fini sur $S$,
$s\in S$, supposons que pour tout $x\in S'$ au dessus de $s$, $X_s$ soit lisse sur
$k(s)$ en $x$ . Alors il existe \unsure{un} voisinage ouvert $U$ de $s$, et un voisinage
ouvert $V$ de $S'|U$ dans $X|U$ , tels que $V\longrightarrow U$ soit poly-élémentaire.
Si $Y$ est un sous-préschéma fermé de $X$\inserted[typed]{ne rencontrant pas $S'$ et} tel que l'ensemble $Z$ des points
en lesquels $Y$ n'est pas lisse sur $S$ \xout{\unsure{aixxdesxfibresxde}} \deleted{soit tel que} \inserted{\unsure{satisfasse}}
dim $Z_s$ \hand{$\leq$} dim $X_s$ - 2 , alors on peut ci-dessus prendre $V$ contenu dans
$X$ - $Y$.

c) (A être éventuellement inclus dans l'énoncé de 13.2.). Soit
avec les notations de 13.2. $\emptyset$ une partie finie de $X$ formée de points
lisses de $X$, et supposons que $\emptyset$ soit contenu dans un \xout{xxx} ouvert affine
de $X$. Alors $\emptyset$ admet un systême fondamental de voisinages ouverts
polyélémentaires. \inserted[typed]{Evidemment on peut supposer $\emptyset$ formé de points \uline{fermés}.} La démonstration est essentiellement la même, sauf
qu'on formule 13.3. sous une forme légèrement différente, savoir: il
existe une sous-variété linéaire $C$ de $P^r$ de codimension $n$, ne rencon\unsure{t}%
trant pas $\emptyset\cup Y$ , et telle que pour tout $x_i\in\emptyset$\inserted{$^{\unsure{(\bar{k})|(k)}}$}, la sous-variété
linéaire $L_i$ de codimension $n$-1 engendrée par $C$ et $x_i$ ait les propriétés
a) , b) , c) de 13.3. Pour vérifier ce point, on note qu'il suffit de
vérifier que la $C$ \uline{générique} a les propriétés envisagées, or \inserted{pour une telle $C$} \deleted{alors}
\xout{\unsure{lxsxlx}}\hand{\unsure{,}} chaque $L_i$ est génér\unsure{ique} parmi les $L$ de codimension $n$ passant
par $x_i$ , de sorte qu'on peut appliquer 13.3. sous la forme initiale
( ou du moins sous la forme où nous l'avions prouvée, qui était: toute
$L_o$ \uline{assez générale} passant par $x$ a les propriétés a),b),c) ).

e) Une des raisons de l'intérêt de 13.2. est la structure``topolo-%
gique'' particulièrement simple des \xout{xxxxxx} \inserted[typed]{schémas alg.} élémentaires\inserted{\unsure{$/U$|$U$}}. Par exemple,

\mspage{\hand{82}}
